概率模型常把答案写成一个期望、一个归一化常数或一个 posterior，但``定义清楚''并不意味着``能够解析计算''。这一单元研究四种容易混淆的随机计算任务：

\begin{enumerate}
\def\labelenumi{\arabic{enumi}.}
\tightlist
\item
  Monte Carlo 用随机样本逼近期望：可以从目标分布独立采样时，用样本平均估计积分，并显式报告随机计算误差。
\item
  MCMC 在无法独立采样时构造平稳链：只能计算未归一化密度时，构造以目标分布为平稳分布的 Markov 链，并处理样本相关性。
\item
  变分推断把后验近似变成优化：不再等待链探索后验，而在易处理的分布族中优化一个近似，以速度换取可分析的近似偏差。
\item
  Bootstrap 重建分析流程的不确定性：它不近似 Bayesian posterior，而是通过模拟重复收集数据，估计整个统计流程在抽样中的波动。
\end{enumerate}

四种方法都产生``许多样本''或``一个分布''，但样本来自哪里、近似哪个对象、误差怎样定义完全不同。读者应能在完成本单元后先识别任务，再选择方法，而不是看到不确定性就机械调用某种采样器。
